\documentclass{article}
\usepackage{lmodern}
\usepackage{amsmath}
\usepackage{listings}
\usepackage{fullpage}
\usepackage{parskip}
\usepackage{xcolor}
\lstset{
	basicstyle=\ttfamily,
	columns=fullflexible,
	frame=single,
	breaklines=true,
	postbreak=\mbox{\textcolor{red}{$\hookrightarrow$}\space},
}
\begin{document}
\title{Experiment `p\_6\_1\_b\_process\_array\_every\_even\_element\_array' Results}
\author{\today}
\date{}
\maketitle
\textbf{Experiment outcome: }SUCCESS
\\\textbf{Bad responses: }0
\\\textbf{Responses containing}~\texttt{assume}~\textbf{: }0
\\\textbf{Responses containing}~\texttt{expect}~\textbf{: }0
\\\textbf{Resolution attempts: }3
\\\textbf{Hard fails (resolution): }1
\\\textbf{Soft fails (resolution): }0
\\\textbf{Verification attempts: }2
\section*{Problem Specification}
\textbf{Problem name: }p\_6\_1\_b\_process\_array\_every\_even\_element\_array
\\\textbf{Natural language statement: }Write a method that takes an array of ten integers and then returns every even element.
\\\textbf{Method signature: }p\_6\_1\_b\_process\_array\_every\_even\_element\_array(inputs: array<int>) returns (even\_elements: array<int>)
\subsection*{Ensures}
\begin{itemize}
\item \texttt{even\_elements[..] == evens(inputs[..])}
\end{itemize}
\subsection*{Requires}
\begin{itemize}
\item \texttt{inputs.Length == 10}
\end{itemize}
\subsection*{Functional Code Given}
\begin{lstlisting}
function evens(inputs: seq<int>): seq<int>
{
  if inputs ==[] then []
  else if inputs[0] % 2 == 0 then [inputs[0]] + evens(inputs[1..])
  else evens(inputs[1..])
}
\end{lstlisting}
\clearpage
\section*{GenAI interactions}
Below you will find all interactions between the `user' (program) and the `assistant' (OpenAI).
\subsection*{Program $\rightarrow$ GenAI}
\begin{lstlisting}
You are given the following task to perform in Dafny:

Write a method that takes an array of ten integers and then returns every even element.

The signature should be:

p_6_1_b_process_array_every_even_element_array(inputs: array<int>) returns (even_elements: array<int>)

The method should respect the following contract:

requires inputs.Length == 10, ensures even_elements[..] == evens(inputs[..])

The contract uses the following dafny code:

function evens(inputs: seq<int>): seq<int>
{
  if inputs ==[] then []
  else if inputs[0] % 2 == 0 then [inputs[0]] + evens(inputs[1..])
  else evens(inputs[1..])
}

These function(s) must not be used in your implementation of the method.

Produce and show only the Dafny body of this method, including the curly braces that surround it. Do not show the signature nor contract. You must not use 'assume' nor 'expect' anywhere in your code.
\end{lstlisting}
\subsection*{GenAI $\rightarrow$ Program}
\textbf{System fingerprint: }null
\\\textbf{ID: }chatcmpl-Dn0p0pBMtqR5oakRZNrpLLjuflyNb
\\\textbf{Model: }gpt-5.5-2026-04-23
\\\textbf{Created at: }1780573602
\\\textbf{Finish reason: }stop
\\\textbf{Usage: }458 tokens in, and 7010 tokens out
\begin{lstlisting}
{
  var s0 := inputs[0..10];
  var s1 := inputs[1..10];
  var s2 := inputs[2..10];
  var s3 := inputs[3..10];
  var s4 := inputs[4..10];
  var s5 := inputs[5..10];
  var s6 := inputs[6..10];
  var s7 := inputs[7..10];
  var s8 := inputs[8..10];
  var s9 := inputs[9..10];
  var s10 := inputs[10..10];

  assert inputs[..] == s0;
  assert |s0| == 10;
  assert |s1| == 9;
  assert |s2| == 8;
  assert |s3| == 7;
  assert |s4| == 6;
  assert |s5| == 5;
  assert |s6| == 4;
  assert |s7| == 3;
  assert |s8| == 2;
  assert |s9| == 1;
  assert |s10| == 0;
  assert s0[1..] == s1;
  assert s1[1..] == s2;
  assert s2[1..] == s3;
  assert s3[1..] == s4;
  assert s4[1..] == s5;
  assert s5[1..] == s6;
  assert s6[1..] == s7;
  assert s7[1..] == s8;
  assert s8[1..] == s9;
  assert s9[1..] == s10;
  assert s10 == [];

  var r10: seq<int> := [];
  var r9: seq<int> := if s9[0] % 2 == 0 then [s9[0]] + r10 else r10;
  var r8: seq<int> := if s8[0] % 2 == 0 then [s8[0]] + r9 else r9;
  var r7: seq<int> := if s7[0] % 2 == 0 then [s7[0]] + r8 else r8;
  var r6: seq<int> := if s6[0] % 2 == 0 then [s6[0]] + r7 else r7;
  var r5: seq<int> := if s5[0] % 2 == 0 then [s5[0]] + r6 else r6;
  var r4: seq<int> := if s4[0] % 2 == 0 then [s4[0]] + r5 else r5;
  var r3: seq<int> := if s3[0] % 2 == 0 then [s3[0]] + r4 else r4;
  var r2: seq<int> := if s2[0] % 2 == 0 then [s2[0]] + r3 else r3;
  var r1: seq<int> := if s1[0] % 2 == 0 then [s1[0]] + r2 else r2;
  var r0: seq<int> := if s0[0] % 2 == 0 then [s0[0]] + r1 else r1;

  even_elements := new int[|r0|](i => r0[i]);
  assert even_elements[..] == r0;
}
\end{lstlisting}
\subsection*{Program $\rightarrow$ GenAI}
\begin{lstlisting}
When using dafny verify, the below error is emitted and verify fails:

problem_data/p_6_1_b_process_array_every_even_element_array/program.dfy(11,0): Error: a postcondition could not be proved on this return path
   |
11 | {
   | ^

problem_data/p_6_1_b_process_array_every_even_element_array/program.dfy(10,27): Related location: this is the postcondition that could not be proved
   |
10 | 	ensures even_elements[..] == evens(inputs[..])
   |                            ^^

problem_data/p_6_1_b_process_array_every_even_element_array/program.dfy(60,40): Error: index out of range
   |
60 |   even_elements := new int[|r0|](i => r0[i]);
   |                                         ^


Dafny program verifier finished with 2 verified, 2 errors

Correct the error in code by altering only the method body. Produce and show only the Dafny body, including the curly braces that surround it. Do not show the signature nor contract. You must not use 'assume' nor 'expect' anywhere in your code.
\end{lstlisting}
\subsection*{GenAI $\rightarrow$ Program}
\textbf{System fingerprint: }null
\\\textbf{ID: }chatcmpl-Dn0s2qtdvkC8trlXGKjGfDWTimpE3
\\\textbf{Model: }gpt-5.5-2026-04-23
\\\textbf{Created at: }1780573790
\\\textbf{Finish reason: }stop
\\\textbf{Usage: }1572 tokens in, and 8639 tokens out
\begin{lstlisting}
{
  var r: seq<int> := [];
  var k := inputs.Length;
  assert k == 10;
  assert inputs[k..] == [];
  assert r == evens(inputs[k..]);

  while k > 0
    invariant inputs.Length == 10
    invariant 0 <= k <= inputs.Length
    invariant r == evens(inputs[k..])
  {
    k := k - 1;
    assert 0 <= k < inputs.Length;
    assert 0 <= k + 1 <= inputs.Length;

    var tail := r;
    assert tail == evens(inputs[k + 1..]);

    var current := inputs[k..];
    assert current != [];
    assert current[0] == inputs[k];
    assert current[1..] == inputs[k + 1..];

    if inputs[k] % 2 == 0 {
      assert current[0] % 2 == 0;
      assert evens(current) == [current[0]] + evens(current[1..]);
      r := [inputs[k]] + tail;
      assert [current[0]] == [inputs[k]];
      assert evens(current[1..]) == tail;
      assert r == evens(current);
    } else {
      assert current[0] % 2 != 0;
      assert evens(current) == evens(current[1..]);
      r := tail;
      assert evens(current[1..]) == tail;
      assert r == evens(current);
    }
    assert current == inputs[k..];
    assert r == evens(inputs[k..]);
  }

  assert k == 0;
  assert inputs[k..] == inputs[..];
  assert r == evens(inputs[..]);

  even_elements := new int[|r|](i => if 0 <= i < |r| then r[i] else 0);
  assert even_elements.Length == |r|;
  assert forall j | 0 <= j < |r|
    ensures even_elements[..][j] == r[j]
  {
    assert even_elements[j] == r[j];
    assert even_elements[..][j] == even_elements[j];
  }
  assert even_elements[..] == r;
  assert even_elements[..] == evens(inputs[..]);
}
\end{lstlisting}
\subsection*{Program $\rightarrow$ GenAI}
\begin{lstlisting}
When using dafny resolve, the below error is emitted and resolve fails:

problem_data/p_6_1_b_process_array_every_even_element_array/program.dfy(67,0): Error: invalid UnaryExpression
   |
67 | }
   | ^

problem_data/p_6_1_b_process_array_every_even_element_array/program.dfy(66,47): Error: expected either 'by' or a semicolon following the assert expression
   |
66 |   assert even_elements[..] == evens(inputs[..]);
   |                                                ^

2 parse errors detected in program.dfy

Correct the error in code by altering only the method body. Produce and show only the Dafny body, including the curly braces that surround it. Do not show the signature nor contract. You must not use 'assume' nor 'expect' anywhere in your code.
\end{lstlisting}
\subsection*{GenAI $\rightarrow$ Program}
\textbf{System fingerprint: }null
\\\textbf{ID: }chatcmpl-Dn0v24jVYmtLka8mJ7W9qnmVrKY8z
\\\textbf{Model: }gpt-5.5-2026-04-23
\\\textbf{Created at: }1780573976
\\\textbf{Finish reason: }stop
\\\textbf{Usage: }2239 tokens in, and 2755 tokens out
\begin{lstlisting}
{
  var r: seq<int> := [];
  var k := inputs.Length;
  assert k == 10;
  assert inputs[k..] == [];
  assert r == evens(inputs[k..]);

  while k > 0
    invariant inputs.Length == 10
    invariant 0 <= k <= inputs.Length
    invariant r == evens(inputs[k..])
  {
    k := k - 1;
    assert 0 <= k < inputs.Length;
    assert 0 <= k + 1 <= inputs.Length;

    var tail := r;
    assert tail == evens(inputs[k + 1..]);

    var current := inputs[k..];
    assert current != [];
    assert current[0] == inputs[k];
    assert current[1..] == inputs[k + 1..];

    if inputs[k] % 2 == 0 {
      assert current[0] % 2 == 0;
      assert evens(current) == [current[0]] + evens(current[1..]);
      r := [inputs[k]] + tail;
      assert [current[0]] == [inputs[k]];
      assert evens(current[1..]) == tail;
      assert r == evens(current);
    } else {
      assert current[0] % 2 != 0;
      assert evens(current) == evens(current[1..]);
      r := tail;
      assert evens(current[1..]) == tail;
      assert r == evens(current);
    }
    assert current == inputs[k..];
    assert r == evens(inputs[k..]);
  }

  assert k == 0;
  assert inputs[k..] == inputs[..];
  assert r == evens(inputs[..]);

  even_elements := new int[|r|](i => if 0 <= i < |r| then r[i] else 0);
  assert even_elements.Length == |r|;

  forall j | 0 <= j < |r|
    ensures even_elements[..][j] == r[j]
  {
    assert even_elements[j] == r[j];
    assert even_elements[..][j] == even_elements[j];
  }

  assert |even_elements[..]| == |r|;
  assert even_elements[..] == r;
  assert even_elements[..] == evens(inputs[..]);
}
\end{lstlisting}
\section*{Final Program}
\begin{lstlisting}
function evens(inputs: seq<int>): seq<int>
{
  if inputs ==[] then []
  else if inputs[0] % 2 == 0 then [inputs[0]] + evens(inputs[1..])
  else evens(inputs[1..])
}

method p_6_1_b_process_array_every_even_element_array(inputs: array<int>) returns (even_elements: array<int>)
	requires inputs.Length == 10
	ensures even_elements[..] == evens(inputs[..])
{
  var r: seq<int> := [];
  var k := inputs.Length;
  assert k == 10;
  assert inputs[k..] == [];
  assert r == evens(inputs[k..]);

  while k > 0
    invariant inputs.Length == 10
    invariant 0 <= k <= inputs.Length
    invariant r == evens(inputs[k..])
  {
    k := k - 1;
    assert 0 <= k < inputs.Length;
    assert 0 <= k + 1 <= inputs.Length;

    var tail := r;
    assert tail == evens(inputs[k + 1..]);

    var current := inputs[k..];
    assert current != [];
    assert current[0] == inputs[k];
    assert current[1..] == inputs[k + 1..];

    if inputs[k] % 2 == 0 {
      assert current[0] % 2 == 0;
      assert evens(current) == [current[0]] + evens(current[1..]);
      r := [inputs[k]] + tail;
      assert [current[0]] == [inputs[k]];
      assert evens(current[1..]) == tail;
      assert r == evens(current);
    } else {
      assert current[0] % 2 != 0;
      assert evens(current) == evens(current[1..]);
      r := tail;
      assert evens(current[1..]) == tail;
      assert r == evens(current);
    }
    assert current == inputs[k..];
    assert r == evens(inputs[k..]);
  }

  assert k == 0;
  assert inputs[k..] == inputs[..];
  assert r == evens(inputs[..]);

  even_elements := new int[|r|](i => if 0 <= i < |r| then r[i] else 0);
  assert even_elements.Length == |r|;

  forall j | 0 <= j < |r|
    ensures even_elements[..][j] == r[j]
  {
    assert even_elements[j] == r[j];
    assert even_elements[..][j] == even_elements[j];
  }

  assert |even_elements[..]| == |r|;
  assert even_elements[..] == r;
  assert even_elements[..] == evens(inputs[..]);
}
\end{lstlisting}
\section*{Total Token Usage}
\textbf{Input tokens: }4269
\\\textbf{Output tokens: }18404
\\\textbf{Reasoning tokens: }16248
\\\textbf{Sum of `total tokens': }22673
\section*{Experiment Timings}
\textbf{Overall Experiment} started at 1780573601760, ended at 1780574054779, lasting 453019ms (453.02 seconds)
\\\textbf{Iteration \#1} started at 1780573601761, ended at 1780573789121, lasting 187360ms (187.36 seconds)
\\\textbf{Iteration \#2} started at 1780573789121, ended at 1780573975402, lasting 186281ms (186.28 seconds)
\\\textbf{Iteration \#3} started at 1780573975490, ended at 1780574054778, lasting 79288ms (79.29 seconds)
\end{document}
